\chapter{Capturer et surveiller le réseau}\label{ch:nw:capturing-and-monitoring-the-network}

Dans la cage simulée de ce chapitre, le journal propre de la stratégie totalise, à la médiane, \qty{1.5}{\micro\second} entre le décodage
d'un paquet et la remise à la \omterm{def:nw:networking-for-trading:frame}{carte réseau} de l'ordre qu'il a provoqué. La capture prélevée aux prises de la cage indique que la trame
de l'ordre quitte le point de remise des ordres \qty{4.4}{\micro\second} après que le premier octet du \omterm{def:nw:capturing-and-monitoring-the-network:trigger}{paquet déclencheur} a atteint le
point de remise des données de marché. L'écart correspond à la portion du chemin que le programme ne voit pas~: la fibre et les
équipements de couche~1, le chemin de réception de la carte vers la mémoire, puis son chemin d'émission. Au 99\up{e}~centile, la capture indique \qty{38}{\micro\second}, et aucun journal écrit par le programme ne pourrait
dire quelle portion du chemin en est responsable.

Ce chapitre traite de l'entreprise qui observe son propre trafic de l'extérieur. Elle copie chaque trame aux points de remise (les
prises du chapitre~2), l'horodate à l'aide des horloges du chapitre~4, la stocke, et répond à partir de la seule capture aux questions que l'entreprise se pose chaque
jour~: combien de temps avons-nous mis, qu'avons-nous perdu, quand le port s'est-il saturé.

\section{Prises, boîtiers de capture et fichiers de capture}

\begin{definition}[Capture de paquets, boîtier de capture]\label{def:nw:capturing-and-monitoring-the-network:capture}
Une \emph{capture de paquets}\index{capture de paquets} est un enregistrement des trames qui ont traversé un point d'un réseau, chacune accompagnée
de l'instant de son passage. Un \emph{boîtier de capture}\index{boîtier de capture} est un équipement dédié qui reçoit des copies du trafic issues des prises,
horodate chaque trame à l'aide d'une horloge disciplinée, et écrit les trames sur un support de stockage au \omterm{def:nw:networking-for-trading:ser}{débit de ligne} de ses entrées sans
en perdre aucune.
\end{definition}

\begin{definition}[Fichier de capture de paquets, suffixe d'horodatage]\label{def:nw:capturing-and-monitoring-the-network:trailer}
Un \emph{fichier de capture de paquets}\index{fichier de capture de paquets} conserve une capture sous la forme d'une suite d'enregistrements, chacun composé
d'un horodatage, d'une longueur et des octets de la trame~; le format libpcap et son successeur pcapng sont les plus répandus, avec des horodatages en nanosecondes dans leurs variantes récentes.
Un \emph{suffixe d'horodatage}\index{suffixe d'horodatage} (timestamp trailer) est un bloc d'octets qu'un équipement de prise ou de commutation ajoute à la fin d'une trame copiée~;
il porte l'instant auquel l'équipement a vu la trame et l'endroit où il l'a vue (équipement et port), de sorte que l'horodatage
voyage avec la trame jusqu'à l'endroit où elle est stockée.
\end{definition}

L'horodatage qui compte est celui pris à l'endroit où la trame a traversé la prise, et non celui que le serveur de stockage écrit lorsque la
copie l'atteint~: entre les deux se trouvent un commutateur d'agrégation, une file d'attente et la carte propre du serveur. Un suffixe y remédie. Il est pris sur le
fil, dans les premiers octets de la trame, et il accompagne la copie à travers toute agrégation ultérieure~; le serveur de stockage le conserve et peut ignorer
sa propre heure. L'\Cref{dat:nw:capturing-and-monitoring-the-network:trailers} décrit une disposition publiée par un constructeur~; le
\texttt{firm.wirecap} du livre utilise un suffixe de 16~octets qui lui est propre, dans le même esprit.

\begin{omfigure}
\begin{tikzpicture}[font=\footnotesize,b/.style={draw,rounded corners=2pt,minimum height=0.6cm,align=center,inner sep=3pt},>=Stealth]
\node[b,fill=gray!10] (A) at (0,1.6) {point A};
\node[b,fill=gray!10] (B) at (0,0.5) {point B};
\node[b,fill=gray!10] (O) at (0,-0.6) {remise des ordres};
\foreach \n/\y in {tA/1.6,tB/0.5,tO/-0.6} {\node[b,fill=omExo!15] (\n) at (2.5,\y) {tap};}
\node[b,fill=omDef!15,minimum height=1.2cm] (cage) at (6.2,0.5) {équipements de la cage\\ et serveurs};
\node[b,fill=omMeth!20,minimum height=1.0cm,minimum width=2.4cm] (agg) at (3.9,-2.6) {horodatage\\ et agrégation};
\node[b,fill=omThm!12,minimum height=1.0cm] (app) at (9.6,-2.6) {boîtier de capture~:\\ stockage, analyse};
\node[b,fill=gray!8] (gm) at (9.6,-0.9) {grand maître\\ (chapitre~4)};
\foreach \a/\b in {A/tA,B/tB,O/tO} {\draw[thick] (\a) -- (\b);}
\foreach \a in {tA,tB,tO} {\draw[thick] (\a.east) -- (cage.west);}
\foreach \a/\x in {tA/3.3,tB/3.55,tO/3.8} {\draw[->,gray,dashed] (\a.south) -- ++(0,-0.25) -| (\x,-2.1);}
\draw[->,thick] (agg) -- node[above,font=\scriptsize]{copies avec suffixes} (app);
\draw[->,gray] (gm) -- node[right,font=\scriptsize]{heure} (app);
\draw[->,gray] (gm.west) -| (4.8,-2.1);
\end{tikzpicture}

\omcaption{Observer une cage de l'extérieur. Des prises sur chaque point de remise copient les deux sens~; un équipement d'horodatage ajoute un suffixe à chaque copie
(instant, équipement, port) par rapport à l'heure du grand maître, puis agrège les copies~; le \omterm{def:nw:capturing-and-monitoring-the-network:capture}{boîtier de capture} les stocke et les analyse. Aucun
programme de la cage n'intervient.}\label{fig:nw:capturing-and-monitoring-the-network:taps}
\end{omfigure}

\begin{omfigure}
\begin{tikzpicture}[font=\scriptsize,f/.style={draw,minimum height=0.8cm,align=center,inner sep=1pt}]
\node[f,fill=omDef!12,minimum width=4.2cm] (fr) at (0,0) {trame d'origine (telle que sur le fil)};
\node[f,fill=omMeth!20,minimum width=1.4cm,anchor=west] (d) at (fr.east) {équipement\\ 2 octets};
\node[f,fill=omMeth!20,minimum width=1.4cm,anchor=west] (p) at (d.east) {port\\ 2 octets};
\node[f,fill=omMeth!20,minimum width=1.4cm,anchor=west] (fl) at (p.east) {indicateurs\\ 4 octets};
\node[f,fill=omThm!20,minimum width=1.5cm,anchor=west] (s) at (fl.east) {secondes\\ 4 octets};
\node[f,fill=omThm!20,minimum width=1.7cm,anchor=west] (n) at (s.east) {nanosecondes\\ 4 octets};
\draw[decorate,decoration={brace,amplitude=4pt,mirror}] (d.south west) -- (n.south east)
  node[midway,below=5pt] {suffixe de 16~octets ajouté par l'équipement de prise (la disposition du livre)};
\end{tikzpicture}

\omcaption{Le \omterm{def:nw:capturing-and-monitoring-the-network:trailer}{suffixe d'horodatage} de \texttt{firm.wirecap}~: l'équipement et le port qui ont vu la trame, ainsi que l'heure UTC complète en secondes et
en nanosecondes, ajoutés après la trame. Les suffixes des constructeurs portent la même information dans leurs propres dispositions.}\label{fig:nw:capturing-and-monitoring-the-network:trailer}
\end{omfigure}

\begin{dated}{2026-09}{Un suffixe publié}\label{dat:nw:capturing-and-monitoring-the-network:trailers}
Arista indique que ses équipements 7130 insèrent un en-tête d'horodatage autonome de 16~octets à la fin de chaque \omterm{def:nw:networking-for-trading:frame}{trame Ethernet}
copiée, recalculent sa séquence de contrôle et portent 64~bits de secondes et de nanosecondes, de sorte que «~l'heure UTC complète~» ne nécessite aucune trame-clé~; l'en-tête
désigne l'équipement et le port, et un indicateur précise si la séquence de contrôle de la trame d'origine était valide. Les horodatages sont pris «~sur le premier
octet d'un paquet au port d'entrée~», le délai de la face avant étant retiré, et le document annonce une précision de $\pm$75~picosecondes lorsque l'équipement est synchronisé
par une impulsion par seconde. D'autres équipements écrasent un champ de la trame par un horodatage plus court et envoient périodiquement
des trames-clés pour le rattacher à l'UTC.
\end{dated}

\omcode{code/firm/wirecap/firm_wirecap.py}{75}{92}{Écriture et lecture d'un fichier pcap en nanosecondes~: un en-tête de fichier de 24~octets, puis pour chaque trame
un en-tête d'enregistrement de 16~octets et les octets.}\label{lst:nw:capturing-and-monitoring-the-network:pcap}

\section{Mesurer la latence wire-to-wire}

La latence wire-to-wire (livre~13, chapitre~1) se mesure sur le fil~: du premier octet d'une trame de données de marché au point
de remise de l'entreprise jusqu'au premier octet de l'ordre qu'elle a provoqué au point de remise des ordres. À partir d'une capture, il s'agit d'une soustraction, une fois que chaque ordre a été
apparié à sa cause.

\begin{definition}[Paquet déclencheur]\label{def:nw:capturing-and-monitoring-the-network:trigger}
Le \emph{paquet déclencheur}\index{paquet déclencheur} d'un ordre est le paquet de données de marché dont le contenu a conduit
l'entreprise à l'envoyer. Son instant de capture marque le début de la latence wire-to-wire de l'ordre.
\end{definition}

L'appariement est la partie difficile. Une capture voit deux flux de paquets, et rien sur le fil n'indique quel paquet a causé quel ordre.
La règle naïve, «~le dernier paquet contenant un message déclencheur avant l'ordre~», fonctionne lorsque les déclencheurs sont rares et que l'entreprise est rapide~;
elle échoue précisément là où la latence vaut la peine d'être mesurée, en rafales, lorsque plusieurs déclencheurs précèdent un ordre de moins que la latence propre de l'entreprise.
La règle fiable consiste à porter la cause dans l'ordre~: l'entreprise place le numéro de séquence du \omterm{def:nw:capturing-and-monitoring-the-network:trigger}{paquet déclencheur} (ou un haché de celui-ci) dans un champ
de l'ordre qu'elle envoie, en général son identifiant d'ordre client, et la capture le relit.

\begin{method}[Wire-to-wire à partir d'une capture]\label{met:nw:capturing-and-monitoring-the-network:w2w}
\begin{enumerate}
\item Prendre des prises sur les deux lignes de données de marché et sur le point de remise des ordres~; horodater les trois à l'aide de la même horloge, sur le fil.
\item Pour chaque paquet de données de marché, conserver l'instant de la première copie vue sur l'une ou l'autre ligne~: la stratégie a agi sur celle qui est arrivée la première.
\item Faire porter à la stratégie le numéro de séquence du déclencheur dans chaque ordre~; pour chaque ordre, retrancher de son propre instant celui de son déclencheur.
\item Présenter la distribution, et non la moyenne, et la ventiler selon le budget de latence des étapes du chemin, afin qu'un centile lent puisse être
  imputé à une étape.
\end{enumerate}
\end{method}

Le tutoriel du chapitre exécute cette méthode sur une seconde simulée~: 99\,055 paquets de données de marché sur deux lignes, chaque copie étant perdue indépendamment
avec une probabilité de 0.2~\%, et une stratégie qui répond par un ordre à 30~\% des paquets portant un message de transaction, soit 9\,558 ordres au total. Le
calendrier des ordres est tiré de \texttt{firm.wirepath}, l'extension au fil du chemin tick-to-trade du livre~13
(\cref{fig:nw:capturing-and-monitoring-the-network:budget}). La capture retrouve exactement la latence de chaque ordre grâce au numéro de séquence qu'il porte,
la règle naïve associe 47~\% d'entre eux au mauvais paquet.

\begin{omfigure}
\begin{tikzpicture}
\begin{axis}[xbar=0pt,width=11.5cm,height=6.6cm,bar width=5pt,xmode=log,xmin=20,xmax=60000,
  ytick=data,yticklabels={réseau entrant,réception carte,décodage,anneau,carnet,stratégie,risque,passerelle,émission carte,réseau sortant,wire-to-wire},
  y dir=reverse,xlabel={nanosecondes (échelle logarithmique)},tick label style={font=\small},label style={font=\small},
  xtick={100,1000,10000},xticklabels={100,1\,000,10\,000},
  grid=major,grid style={gray!25},legend columns=2,legend style={font=\footnotesize,at={(0.5,-0.2)},anchor=north,draw=none,
  fill=none,/tikz/every even column/.append style={column sep=8pt}},table/col sep=comma]
\addplot[fill=omDef!55,draw=omDef] table[y expr=\coordindex,x=p50_ns]{figdata/networks/05-capturing-and-monitoring-the-network/budget.csv};
\addplot[fill=omThm!40,draw=omThm] table[y expr=\coordindex,x=p99_ns]{figdata/networks/05-capturing-and-monitoring-the-network/budget.csv};
\legend{médiane,centile 99}
\end{axis}
\end{tikzpicture}

\omcaption{Budget fil à fil de \texttt{firm.wirepath}~: la cage de couche~1 du chapitre~2 (réseau entrant et sortant), les chemins de réception et
d'émission de la carte (valeurs de modèle), et les six étapes internes au processus mesurées par le livre~13 sur un ordinateur portable (Intel Core Ultra 7 155H, WSL2).
La dernière barre est leur composition, une simulation. Données~: \texttt{fig\_capture.py}.}\label{fig:nw:capturing-and-monitoring-the-network:budget}
\end{omfigure}

Le budget indique où part le temps. À la médiane, les deux passages de la carte (\qty{0.8}{\micro\second} chacun, valeurs de modèle) et l'anneau
entre le feed handler et la stratégie (\qty{0.76}{\micro\second}, mesuré) sont les étapes les plus longues~; la fibre et les équipements de la cage
totalisent ensemble \qty{0.56}{\micro\second}. Au centile 99, c'est l'anneau qui domine, à \qty{32}{\micro\second}~: sur l'ordinateur portable où il a été mesuré,
un consommateur qui s'endort ou qui est désordonnancé se réveille en retard. Composées, les étapes donnent une latence médiane fil à fil de \qty{4.4}{\micro\second} et un centile 99
de \qty{35}{\micro\second}~; la capture de la session simulée mesure \qty{4.4}{\micro\second} et \qty{38}{\micro\second}, la différence
tenant à l'échantillon.

\begin{omfigure}
\begin{tikzpicture}
\begin{axis}[width=11.5cm,height=5.0cm,xmode=log,xmin=2,xmax=100,ymin=0,ymax=0.16,
  xlabel={latence fil à fil issue de la capture (\unit{\micro\second}, échelle logarithmique)},ylabel={part des ordres},
  ytick={0,0.05,0.1,0.15},yticklabels={0,0.05,0.10,0.15},
  xtick={2,5,10,20,50,100},xticklabels={2,5,10,20,50,100},
  tick label style={font=\small},label style={font=\small},grid=major,grid style={gray!25},table/col sep=comma]
\addplot[ybar,bar width=3pt,fill=omDef!55,draw=omDef] table[x=mid_us,y=share]{figdata/networks/05-capturing-and-monitoring-the-network/w2w_hist.csv};
\end{axis}
\end{tikzpicture}

\omcaption{Simulation~: latence fil à fil de 9\,558 ordres mesurée à partir de la capture d'une seconde simulée, selon la séquence de déclenchement
que porte chaque ordre (classes de même largeur en échelle logarithmique). La distribution présente la longue queue droite de son étape la plus lente. Données~:
\texttt{fig\_capture.py}.}\label{fig:nw:capturing-and-monitoring-the-network:hist}
\end{omfigure}

\omcode{code/networks/05-capturing-and-monitoring-the-network/python/nw_capture.py}{74}{95}{Lecture de la capture~: le temps de chaque trame
issu de son suffixe, la première copie de chaque paquet de données de marché sur l'une ou l'autre ligne, et la séquence de déclenchement de chaque
ordre.}\label{lst:nw:capturing-and-monitoring-the-network:analyse}

\section{Détecter les trous, les pertes et les micro-rafales à partir d'une capture}

La même capture répond aux questions que le feed handler du livre~13 traite de l'intérieur, et à certaines qu'il ne peut pas traiter. Un trou est un numéro de
séquence manquant sur une ligne~; un trou sur les deux lignes correspond à un message que la firme n'a jamais reçu, et la capture indique s'il a été perdu
avant la remise (absent à la prise) ou après (présent à la prise, absent des compteurs du handler). Dans la seconde simulée,
la ligne~A présente 158 trous et la ligne~B 198, et aucun message n'est absent des deux lignes~: le handler a tout récupéré sur l'autre ligne. Les rafales
se lisent dans les mêmes horodatages~: les 100 microsecondes les plus chargées de la ligne~A ont transporté \qty{1.24}{\giga\bit\per\second}, soit neuf fois sa moyenne de
\qty{0.14}{\giga\bit\per\second}, la même forme que les flux du chapitre~1 à plus petite échelle.

Une capture est la preuve de la firme autant que son instrument. Elle montre ce que la firme a effectivement envoyé et quand, ce sur quoi une plateforme, un client ou un
régulateur interroge après un incident~; les règles européennes demandent aux firmes d'identifier le point exact où chaque horodatage est appliqué, et une prise sur la
remise est le seul point qu'aucune modification logicielle ne peut déplacer.

\section{Surveiller en production~: compteurs, télémétrie et alarmes}

\begin{definition}[Télémétrie en flux]\label{def:nw:capturing-and-monitoring-the-network:telemetry}
La \emph{télémétrie en flux}\index{télémétrie en flux} (streaming telemetry) est l'export continu, par les équipements réseau, de leurs compteurs et de leur état
(trafic par port, rejets, occupation des files, niveaux optiques, état de l'horloge) vers un collecteur, poussé à un intervalle choisi ou à chaque changement,
au lieu d'être interrogé par le collecteur.
\end{definition}

Une capture est détaillée et coûteuse~: elle conserve chaque trame, et une cage chargée remplit des disques en quelques heures. La surveillance fonctionne en continu
sur des synthèses~: les compteurs de chaque équipement, diffusés chaque seconde ou à chaque changement~; les rejets de socket et d'anneau de chaque serveur~; les compteurs
de trous et d'obsolescence de chaque handler~; le décalage de chaque horloge. Les alarmes qui comptent dans un réseau de trading sont peu nombreuses et précises~:
\begin{itemize}
\item \textbf{rejets sur une \omterm{def:nw:networking-for-trading:egress}{file de sortie}}, que seules les \omterm{def:nw:networking-for-trading:burst}{micro-rafales} provoquent sur une liaison sous-utilisée (chapitre~1)~;
\item \textbf{trous sur une ligne} qui persistent, une ligne qui se dégrade pendant que l'autre la masque~;
\item \textbf{trous sur les deux lignes}, cause commune, et toute reprise par instantané (snapshot)~;
\item \textbf{une horloge hors verrouillage ou en maintien}, et des décalages qui approchent le budget (chapitre~4)~;
\item \textbf{centiles fil à fil} qui s'écartent de la référence, calculés à partir de la capture toutes les minutes.
\end{itemize}
Chaque alarme désigne un point du chemin, et la capture des minutes qui l'entourent est conservée pour l'enquête.

\section{Tutoriel~: une seconde de la cage, depuis les prises}\label{tut:nw:capturing-and-monitoring-the-network:tut}

\begin{tutorial}
\textbf{Objectif.} Capturer une seconde simulée aux prises de la cage, l'écrire sous forme de fichier pcap, et mesurer à partir de lui seul ce que
les programmes de la firme ne peuvent pas voir. \textbf{État final~:} \cref{fig:nw:capturing-and-monitoring-the-network:budget,fig:nw:capturing-and-monitoring-the-network:hist},
l'erreur d'appariement de la règle naïve, et les comptes de trous et de rafales cités dans le texte.
\begin{tutsteps}
\item \textbf{Le budget.} \texttt{nw\_capture.stages()} prend la cage du chapitre~2, les étapes de modèle de la carte et les étapes mesurées du livre~13~;
  \texttt{budget\_rows()} les compose avec \texttt{firm.wirepath}.
\item \textbf{La session.} \texttt{session()} construit les trames du flux sur les deux lignes avec des pertes indépendantes, les ordres avec leur
  séquence de déclenchement, et ajoute un suffixe à chaque copie.
\item \textbf{Le fichier.} \texttt{firm.wirecap.write\_pcap(path, records)} l'écrit (\cref{lst:nw:capturing-and-monitoring-the-network:pcap})~;
  tout outil qui lit les fichiers pcap à la nanoseconde peut l'ouvrir.
\item \textbf{L'analyse.} \texttt{analyse(read\_pcap(path))} apparie les ordres à leurs déclencheurs, compte les trous par ligne et les messages perdus sur
  les deux lignes, et repère la fenêtre la plus chargée (\cref{lst:nw:capturing-and-monitoring-the-network:analyse})~; \texttt{python fig\_capture.py} écrit les
  données des graphiques.
\end{tutsteps}
\textbf{Pour aller plus loin.} Porter la perte à 5~\% par ligne et compter les messages perdus sur les deux lignes~; rendre la stratégie quatre fois plus rapide et
mesurer comment l'erreur de la règle naïve évolue.
\end{tutorial}

\section{Construction~: outils de capture et chemin fil à fil}\label{bld:nw:capturing-and-monitoring-the-network:wirecap}

\begin{build}
\bfield{Rôle} Deux composants~: \texttt{firm.wirecap}, la boîte à outils de capture de la firme, et \texttt{firm.wirepath}, le chemin tick-to-trade du
livre~13 étendu au fil, que les chapitres~7 et~29 utilisent pour leurs tableaux de latence.
\bfield{Interface} \texttt{firm\_wirecap}~: \texttt{udp\_frame}, \texttt{tcp\_frame}, \texttt{add\_trailer}, \texttt{strip\_trailer},
\texttt{write\_pcap}, \texttt{read\_pcap}, \texttt{write\_pcapng}, \texttt{read\_pcapng}, \texttt{parse}, \texttt{mold}, \texttt{ORDER},
\texttt{match\_orders}, \texttt{match\_nearest}, \texttt{gaps}, \texttt{bursts}~; lecteur C++20 \texttt{cpp/firm\_wirecap.hpp}.
\texttt{firm\_wirepath}~: \texttt{Stage(name, p50, p99, kind)}, \texttt{book13\_software()}, \texttt{wire\_stages(cage, \dots)},
\texttt{card\_stages}, \texttt{hardware\_stage(cycles, mhz)}, \texttt{compose}, \texttt{report}, composé avec \texttt{firm.latbudget}.
\bfield{Règles} Les horodatages sont des entiers de nanosecondes depuis minuit UTC~; l'heure du suffixe l'emporte sur toute heure de stockage~; la somme de contrôle IPv4 est
correcte~; un ordre n'est apparié que par la séquence qu'il porte~; les mesures du livre~13 sont lues, jamais recopiées ni modifiées.
\bfield{Tests d'acceptation} \texttt{code/firm/wirecap/tests/}~: les trames se relisent à l'identique, les sommes de contrôle tiennent, les suffixes font l'aller-retour, les fichiers pcap et pcapng
font l'aller-retour, un mauvais numéro magique est refusé, l'appariement, les trous et les rafales sur des données fabriquées à la main, et le jeu de données partagé se lit de la même façon en Python et
en C++. \texttt{code/firm/wirepath/tests/}~: l'ajustement log-normal retrouve sa médiane et son centile 99, les étapes du livre~13 sont lues, le chemin d'une cage
donne l'étape réseau, et les lignes du rapport.
\bfield{Pour aller plus loin} Détection des rafales par port, avec le tampon du commutateur comme seuil~; reconstruction du carnet d'ordres à partir de la capture, comparée,
message par message, à celle du handler.
\end{build}

\begin{omsources}
\item \texttt{pcap-savefile(5)} (tcpdump/libpcap)~; IETF, \emph{PCAP Now Generic (pcapng) Capture File Format},
draft-ietf-opsawg-pcapng-04.
\item Arista, \emph{An Overview of Arista Ethernet Capture Timestamps}~; OpenConfig, \emph{gNMI specification}.
\item Commission Delegated Regulation (EU) 2017/574 (RTS~25), article 4.
\end{omsources}

\section{Exercices}

\begin{exercise}[$\star$]\label{exo:nw:capturing-and-monitoring-the-network:1}
Un en-tête d'enregistrement pcap en nanosecondes donne 34\,200 secondes et 1\,500\,000\,321 nanosecondes. Vérifier si cet enregistrement est valide,
et déterminer l'heure de la journée que doit indiquer le lecteur.
\end{exercise}

\begin{exercise}[$\star$]\label{exo:nw:capturing-and-monitoring-the-network:2}
Un \omterm{def:nw:capturing-and-monitoring-the-network:trigger}{paquet déclencheur} atteint le tap sur la ligne~A à 09:30:00.000\,012\,345 et sur la ligne~B \qty{3}{\micro\second} plus tard~; l'ordre part à
09:30:00.000\,016\,771. Calculer la latence wire-to-wire. Si la copie de la ligne~A avait été perdue, à partir de quelle copie faudrait-il mesurer
et préciser ce qui changerait.
\end{exercise}

\begin{exercise}[$\star$]\label{exo:nw:capturing-and-monitoring-the-network:3}
Une liaison à 10~Gb/s est tappée dans les deux sens et écrite sur disque avec ses \omterm{def:nw:capturing-and-monitoring-the-network:trailer}{suffixes d'horodatage}~; chaque trame est stockée avec un en-tête d'enregistrement de 16~octets et
son suffixe de 16~octets. Calculer, à un débit soutenu de \qty{4}{\giga\bit\per\second} dans chaque sens et avec des trames de 200~octets, le nombre d'octets par seconde écrits sur disque.
\end{exercise}

\begin{exercise}[$\star\star$]\label{exo:nw:capturing-and-monitoring-the-network:4}
Expliquer pourquoi la règle naïve d'appariement échoue d'autant plus souvent que les déclencheurs sont fréquents et que la firme est lente.
\end{exercise}

\begin{exercise}[$\star\star$]\label{exo:nw:capturing-and-monitoring-the-network:5}
La ligne~A présente 158 lacunes et la ligne~B 198 dans la seconde du tutoriel, et aucune n'est commune aux deux. Avec une perte indépendante de 0.2~\% par copie et 99\,055
paquets, calculer le nombre de paquets qu'on attend de perdre sur les deux lignes.
\end{exercise}

\begin{exercise}[$\star\star$]\label{exo:nw:capturing-and-monitoring-the-network:6}
Dans la \cref{fig:nw:capturing-and-monitoring-the-network:budget}, indiquer quelle étape attaquer en premier pour réduire la médiane, et laquelle pour réduire le
99\up{e}~centile. Préciser le coût de chacune.
\end{exercise}

\begin{exercise}[$\star\star\star$]\label{exo:nw:capturing-and-monitoring-the-network:7}
\emph{Programmation.} Rendre chaque étape de \texttt{firm.wirepath} quatre fois plus rapide, sauf celle du réseau, relancer la session et mesurer
l'erreur d'appariement de la règle naïve, puis expliquer le sens de la variation.
\end{exercise}

\begin{exercise}[$\star\star\star$]\label{exo:nw:capturing-and-monitoring-the-network:8}
\emph{Trouver la faille.} «~Nous mesurons notre latence wire-to-wire en horodatant le paquet de données de marché dans le handler avec
l'horodatage matériel de la carte, et l'ordre avec l'horodatage d'émission de la carte. C'est aussi bon qu'un tap.~»
\end{exercise}

\section{Problème~: Les deux microsecondes que personne n'a journalisées}

\begin{problem}[{Problème du week-end --- un budget de latence lu sur les taps}]\label{pb:nw:capturing-and-monitoring-the-network:1}
La stratégie de la firme journalise son temps de traitement interne et annonce une médiane de \qty{1.5}{\micro\second}. La capture aux taps (la
seconde simulée du chapitre) donne une médiane wire-to-wire de \qty{4.4}{\micro\second} et un 99\up{e}~centile de \qty{38}{\micro\second}. Utiliser les
médianes et les 99\up{e}~centiles par étape de la \cref{fig:nw:capturing-and-monitoring-the-network:budget}~: réseau entrant 263 et sortant 295~ns~; réception carte 798
(p99 1\,997) et émission 801 (p99 2\,000)~; les six étapes logicielles 437, 756, 158, 97, 47 et 33~ns à la médiane.

\medskip\noindent\textbf{Partie~I --- Médianes.}
\begin{enumerate}
\item Calculer la somme des médianes des étapes logicielles.
\item Calculer la somme des médianes de toutes les étapes.
\item Expliquer pourquoi la somme des médianes n'est pas la médiane de la somme.
\item Déterminer la part de la médiane wire-to-wire située hors du programme.
\end{enumerate}

\medskip\noindent\textbf{Partie~II --- La queue de distribution.}
\begin{enumerate}[resume]
\item Le 99\up{e}~centile du ring est de \qty{31.9}{\micro\second}. Déterminer la part qu'il représente dans le 99\up{e}~centile wire-to-wire.
\item Expliquer pourquoi une seule étape lente peut dominer la queue de distribution.
\item Indiquer ce que le livre~13 a trouvé pour la queue du ring sur l'ordinateur portable.
\item Indiquer ce qu'une machine de production modifierait.
\end{enumerate}

\medskip\noindent\textbf{Partie~III --- La capture.}
\begin{enumerate}[resume]
\item Expliquer pourquoi l'horodatage du tap est le bon point de départ de la latence, plutôt que celui de la carte.
\item La règle naïve a mal apparié 47~\% des ordres. Indiquer à quoi aurait ressemblé le 99\up{e}~centile avec elle.
\item Indiquer le champ à utiliser pour transporter le déclencheur dans l'ordre.
\item Indiquer quelle alarme aurait signalé à la firme que sa queue avait bougé.
\end{enumerate}

\medskip\noindent\textbf{Partie~IV --- Le verdict.}
\begin{enumerate}[resume]
\item Énoncer le \emph{résultat nommé}~: la médiane et le 99\up{e}~centile wire-to-wire, la part hors du programme et l'étape qui domine la
  queue.
\item Calculer de combien baisserait la médiane si les chemins des cartes étaient divisés par deux.
\item Déterminer quelle étape dominerait le 99\up{e}~centile si la queue du ring était supprimée entièrement.
\item Indiquer ce qu'il faudrait mesurer pour évaluer les étapes des cartes au lieu de les supposer.
\item Expliquer pourquoi la capture fait partie des preuves de la firme autant que de ses outils.
\item Calculer l'espace disque occupé par une heure de cette capture, aux débits du tutoriel.
\item Indiquer ce qu'il faudrait garder un mois, et ce qu'il faudrait garder un jour.
\item En une phrase~: dire ce que voit un tap et que ne voit pas un journal.
\end{enumerate}
\end{problem}

\section{Questions d'entretien}

\begin{interviewq}[$\star$ \iqroles{développeur}]\label{iq:nw:capturing-and-monitoring-the-network:1}
Comment mesureriez-vous la latence wire-to-wire~? De quoi auriez-vous besoin en plus des journaux propres au programme~?
\end{interviewq}

\begin{interviewq}[$\star\star$ \iqroles{développeur}]\label{iq:nw:capturing-and-monitoring-the-network:2}
Où prendriez-vous l'horodatage d'une capture, et pourquoi cela a-t-il de l'importance~?
\end{interviewq}

\begin{interviewq}[$\star\star$ \iqroles{développeur, chercheur}]\label{iq:nw:capturing-and-monitoring-the-network:3}
Comment appariez-vous chaque ordre au paquet de données de marché qui l'a déclenché~?
\end{interviewq}

\begin{interviewq}[$\star\star$ \iqroles{développeur}]\label{iq:nw:capturing-and-monitoring-the-network:4}
Quelles alarmes réseau mettriez-vous en place pour une cage de trading, et pourquoi celles-là~?
\end{interviewq}

\begin{interviewq}[$\star\star\star$ \iqroles{développeur}]\label{iq:nw:capturing-and-monitoring-the-network:5}
Votre 99\up{e}~centile wire-to-wire a doublé pendant la nuit, alors qu'aucun déploiement n'a eu lieu. Comment en trouveriez-vous la cause~?
\end{interviewq}

\begin{interviewq}[$\star\star$ \iqroles{développeur, chercheur}]\label{iq:nw:capturing-and-monitoring-the-network:6}
À quoi utiliseriez-vous une capture, en dehors de la latence~?
\end{interviewq}
